% A turn of a T-mesh, its corridor, its shift and the helix it makes.
% \input inside a figure; needs \usepackage{tikz} and
% \usetikzlibrary{arrows.meta,decorations.pathreplacing}.
%
% (a) A band of a T-mesh on a cylinder, cut open along the left side of P_1:
% the right side of P_4 is glued to the left side of P_1, drawn again, dashed,
% where the walk comes back into it. The band continues above and below. Every
% node is at a whole position; the faint grid is the quad mesh at these
% lengths. The turn gamma starts on the left side of P_1 at offset o = 2.5 and
% crosses P_1, P_2, P_3, P_4. Measured from the first (lower) corner of each
% patch, its offset grows by l_b crossing into P_2, drops by l_c into P_3 and
% grows by l_d into P_4; the bottoms of P_4 and P_1 meet at the cut. So it comes
% back into P_1 at o + s with s = l_b - l_c + l_d = 1 - 2 + 2 = 1. The corridor
% (the offsets whose walk crosses the same arcs) is 1 < o < 4.
% (b) The strip of quads at o = 2.5 on the cylinder: each round comes back one
% row beside itself, at o = 1.5, 2.5, 3.5 within the corridor, then leaves it:
% a helix of pitch 1. The numbers are the offsets where the rounds start.
\begin{tikzpicture}[x=0.56cm, y=0.56cm, font=\small,
    side/.style={draw=black, line width=0.6pt},
    ghost/.style={draw=black!55, line width=0.6pt, dashed},
    quadgrid/.style={draw=black!18, line width=0.3pt},
    node dot/.style={circle, fill=black, inner sep=0pt, minimum size=3.2pt},
    arc mark/.style={line width=2.2pt, draw=black},
    turn/.style={line width=1.1pt, draw=black, -{Stealth[length=6pt]}},
    brace/.style={decorate, decoration={brace, amplitude=4pt}, line width=0.6pt}]
  \definecolor{corridor}{HTML}{CFE0F3}
  \definecolor{loopspiral}{HTML}{D55E00}

  % ---------------------------------------------------------------- (a)
  \begin{scope}
    \clip (-0.1,-2.4) rectangle (16.1,7.4);
    % The corridor of the turn, 1 < o < 4, from P_1 round to P_1 again.
    \fill[corridor] (0,1) rectangle (13,4);
    \draw[quadgrid, step={(1,1)}] (0,-3) grid (16,8);
    % Sides of the columns of patches, which run through the band.
    \foreach \x in {0,3,7,10,13} \draw[side] (\x,-3) -- (\x,8);
    \draw[ghost] (16,-3) -- (16,8);
    % The bottom and top of P_1, ..., P_4, and of P_1 again.
    \foreach \x/\X/\b/\t in {0/3/0/5, 3/7/-1/4, 7/10/1/6, 10/13/-1/4} {
      \draw[side] (\x,\b) -- (\X,\b) (\x,\t) -- (\X,\t);
      \foreach \px in {\x,\X} \foreach \py in {\b,\t}
        \node[node dot] at (\px,\py) {};
    }
    \draw[ghost] (13,-1) -- (16,-1) (13,4) -- (16,4);
    \foreach \py in {-1,4} \node[node dot, fill=black!55] at (16,\py) {};
  \end{scope}
  % The arcs whose lengths the shift adds up.
  \draw[arc mark] (3,-1) -- (3,0) node[midway, left=1pt] {$b$};
  \draw[arc mark] (7,-1) -- (7,1) node[midway, left=1pt] {$c$};
  \draw[arc mark] (10,-1) -- (10,1) node[midway, right=1pt] {$d$};
  % Patch names.
  \node at (1.5,4.5) {$P_1$};
  \node at (5,0) {$P_2$};
  \node at (8.5,5) {$P_3$};
  \node at (11.5,0) {$P_4$};
  \node[black!60] at (14.5,3.35) {\footnotesize $P_1$ again};
  % The turn, from offset o = 2.5 of the first side of P_1.
  \draw[turn] (0,2.5) -- (13.08,2.5);
  \node[above] at (6.4,2.5) {$\gamma$};
  \draw[black, line width=0.6pt, fill=white] (0,2.5) circle (2.2pt);
  % Where it comes back: in the frame of P_1, at o + s.
  \draw[black!60, dashed, line width=0.8pt] (13,2.5) -- (16,2.5);
  \draw[black!60, dotted, line width=0.9pt] (13,1.5) -- (16,1.5);
  \draw[black, line width=0.6pt, fill=white] (13,1.5) circle (2.2pt);
  \node[right, inner sep=1pt] at (16.15,2.5) {\footnotesize $o + s_\gamma(\ell)$};
  \node[right, inner sep=1pt] at (16.15,1.5) {\footnotesize $o$};
  \draw[brace, decoration={mirror}] (15.7,1.5) -- (15.7,2.5)
    node[midway, left=4pt, inner sep=1pt] {\footnotesize $s_\gamma$};
  % Offset and corridor on the first side.
  \node[left] at (-0.05,2.5) {$o$};
  \draw[brace] (-0.9,1) -- (-0.9,4) node[midway, left=4pt, rotate=90, anchor=south]
    {\footnotesize corridor};
  % The shift as a linear form of the arc lengths.
  \node[anchor=north west] at (0,-2.7)
    {$s_\gamma(\ell) = \ell_b - \ell_c + \ell_d = 1 - 2 + 2 = 1$};
  \node[anchor=north] at (8,-4.2) {(a)};

  % ---------------------------------------------------------------- (b)
  % The strip of quads at o = 2.5 on the cylinder, seen from slightly above:
  % a helix of pitch |s| = 1 row, three rounds in the corridor, 13 quads a
  % round (the length of the turn).
  \begin{scope}[shift={(22.4,-1.6)}]
    \pgfmathsetmacro{\R}{1.9}   % radius
    \pgfmathsetmacro{\k}{0.2}   % flattening of the ellipses
    \pgfmathsetmacro{\w}{1.15}  % one row of quads
    \pgfmathsetmacro{\H}{7.0}   % height
    \pgfmathsetmacro{\zz}{1.4}  % bottom of the first round
    % Back halves: the centre line of the strip, seen through the cylinder.
    \foreach \j in {0,1,2} {
      \draw[loopspiral!80, dashed, line width=0.7pt]
        plot[domain=360*\j+180:360*\j+360, samples=40, variable=\t]
          ({\R*cos(\t)}, {\zz + \w*(\t/360 + 0.5) - \k*\R*sin(\t)});
    }
    % The cylinder, translucent.
    \fill[black!25, fill opacity=0.18] (-\R,0) -- (-\R,\H) arc[start angle=180, end angle=360,
      x radius=\R, y radius=\k*\R] -- (\R,0) arc[start angle=0, end angle=-180,
      x radius=\R, y radius=\k*\R] -- cycle;
    \draw[line width=0.6pt] (-\R,0) -- (-\R,\H) (\R,0) -- (\R,\H);
    \draw[line width=0.6pt] (0,\H) ellipse[x radius=\R, y radius=\k*\R];
    \draw[line width=0.6pt] (\R,0) arc[start angle=0, end angle=-180,
      x radius=\R, y radius=\k*\R];
    \draw[line width=0.4pt, black!40, dashed] (\R,0) arc[start angle=0, end angle=180,
      x radius=\R, y radius=\k*\R];
    % Front halves: the quads of the strip, 13 a round, and its centre line.
    \foreach \j in {0,1,2} {
      \fill[loopspiral!35, fill opacity=0.55]
        plot[domain=360*\j:360*\j+180, samples=40, variable=\t]
          ({\R*cos(\t)}, {\zz + \w*\t/360 - \k*\R*sin(\t)})
        -- plot[domain=360*\j+180:360*\j, samples=40, variable=\t]
          ({\R*cos(\t)}, {\zz + \w*(\t/360 + 1) - \k*\R*sin(\t)}) -- cycle;
      \foreach \e in {0,1} {
        \draw[loopspiral!70, line width=0.5pt]
          plot[domain=360*\j:360*\j+180, samples=40, variable=\t]
            ({\R*cos(\t)}, {\zz + \w*(\t/360 + \e) - \k*\R*sin(\t)});
      }
      \foreach \q in {1,...,6} {
        \pgfmathsetmacro{\t}{360*\j + \q*360/13}
        \draw[loopspiral!70, line width=0.4pt]
          ({\R*cos(\t)}, {\zz + \w*\t/360 - \k*\R*sin(\t)}) --
          ({\R*cos(\t)}, {\zz + \w*(\t/360 + 1) - \k*\R*sin(\t)});
      }
      \draw[loopspiral, line width=1.1pt, -{Stealth[length=5pt]}]
        plot[domain=360*\j:360*\j+176, samples=40, variable=\t]
          ({\R*cos(\t)}, {\zz + \w*(\t/360 + 0.5) - \k*\R*sin(\t)});
      % The offset on the first side where the round starts.
      \pgfmathsetmacro{\n}{\j + 1.5}
      \node[right, font=\scriptsize, inner sep=2pt] at (\R,\zz + \w*\j + 0.5*\w)
        {$\pgfmathprintnumber[fixed, precision=1, fixed zerofill]{\n}$};
    }
    \node[anchor=north, align=center, font=\footnotesize] at (0,-0.5)
      {the strip at $o = 2.5$: one helix\\ through $o = 1.5$, $2.5$, $3.5$,\\ pitch $|s_\gamma(\ell)| = 1$};
  \end{scope}
  \node[anchor=north] at (22.4,-4.2) {(b)};
\end{tikzpicture}
